\chapter{Ensamble de la placa}\label{cap:ensamble}

Herramientas: cautín y soldadura, multímetro, destornillador pequeño (para el potenciómetro del LM2596 y las borneras), fuente de 12~V y lupa.

\section{Antes de empezar}

\begin{enumerate}
\item Revisa la placa desnuda: sin pistas cortadas ni puentes de soldadura.
\item Si es un lote nuevo, compara la placa ESP32-S3 real con la huella: 2 hileras de 22 pines, paso 2,54~mm, 22,86~mm entre hileras, y la posición de 5V, GND y los pines de datos. Los dos conectores USB-C deben quedar accesibles.
\end{enumerate}

\section{Soldar los componentes}

Suelda de lo más bajo a lo más alto:
\begin{enumerate}
\item R1 y R2.
\item C1 y C2, respetando la polaridad (la franja del condensador es el negativo).
\item Portafusibles F1, F2, F5, F6 y F7.
\item Borneras: S\_LM\_INY1, S\_A\_B1, INY1, INY2, A1, B1 y B2, con la entrada de cable hacia el borde de la placa.
\item Módulo LM2596.
\end{enumerate}
Resultado esperado: todos los componentes menos la ESP32 en su lugar.

\section{Ajustar el regulador a \dato{reguladorSalida}}

\advertencia{Ajusta el LM2596 \textbf{antes} de poner la ESP32. Muchos módulos vienen de fábrica con una salida mayor de 5~V, que daña la ESP32 al primer encendido.}

\begin{enumerate}
\item Conecta 12~V a S\_LM\_INY1 (borne 2 positivo, borne 1 GND).
\item Mide con el multímetro entre OUT+ y OUT- del módulo.
\item Gira el potenciómetro azul hasta leer \dato{reguladorSalida}.
\item Comprueba la polaridad: OUT+ debe llegar al pin 5V de la huella de la ESP32 y OUT- a GND.
\item Desconecta los 12~V.
\end{enumerate}

\foto{lm2596-ajuste}{Multímetro midiendo 5,0 V en la salida del LM2596 mientras se gira el potenciómetro.}{0.7}

\section{Instalar la ESP32-S3}

\begin{enumerate}
\item Coloca el módulo con la antena hacia la zona sin cobre de la placa y los USB-C accesibles.
\item Suéldalo (o insértalo en zócalo, si la placa los lleva).
\item Sin conectar el USB, aplica 12~V en S\_LM\_INY1 y mide 5~V entre el pin 5V del módulo y GND.
\end{enumerate}

\advertencia{No conectes al mismo tiempo el USB de la ESP32 y la alimentación del LM2596. Espressif indica que alimentar por USB y por el pin 5V son opciones excluyentes. Para grabar o depurar por USB, desconecta primero los 12~V de la placa.}

\foto{esp-instalada}{ESP32-S3 instalada en la placa, con la antena y los dos USB-C visibles.}{0.7}

\section{Comprobaciones finales sin carga}

Con la placa desconectada, mide continuidad y confirma que \textbf{no} hay corto entre:
\begin{itemize}
\item +12~V y GND en S\_LM\_INY1, S\_A\_B1, A1, B1, INY1 e INY2.
\item 5V y GND de la ESP32.
\end{itemize}
Pon los fusibles (capítulo~\ref{cap:cableado}). Anota los valores medidos en la hoja de verificación (anexo~D).
